\chapter{\MakeUppercase{Background}}
\setcounter{equation}{0}

\section{{Stroke and Motor Impairment}}

Stroke can damage brain regions responsible for controlling voluntary movement. Depending on the
location and severity of the damage, patients may experience weakness, abnormal muscle tone,
reduced joint movement and difficulty maintaining balance. Lower-limb impairment can affect walking,
standing and the ability to perform daily activities independently.
Recovery is influenced by several factors, including initial motor impairment, muscle spasticity, balance
function, age and time elapsed since stroke. These factors make individual outcome prediction a
complex problem.

\section{{Brain–Computer Interface}}

A Brain–Computer Interface establishes a communication pathway between brain activity and an
external device without relying entirely on conventional muscle movement. In the selected study, a BCI
system for lower-limb rehabilitation (BCI-LE) combines electroencephalogram (EEG) signal acquisition
with a robotic ankle-stretching device.
\begin{itemize}
    \item EEG electrodes record the patient's brain activity.
    \item The patient is instructed to imagine moving the affected ankle.
    \item Signal-processing techniques extract relevant features from EEG signals.
    \item A trained classifier identifies the intended motor imagery.
    \item When the intended movement is recognized, the robotic device assists ankle movement.
    \item The patient receives visual and proprioceptive feedback during the process.
\end{itemize}
The objective is to associate intended movement with assisted physical movement. This may support
rehabilitation through repeated, feedback-based training. However, patients can differ in their response
to BCI therapy, making outcome prediction an important research problem.\cite{bib3}


\section{Machine Learning in Rehabilitation}
Machine learning enables computer systems to learn relationships from data and use these patterns to
make predictions. In stroke rehabilitation, patient information can include demographic details, clinical
assessment scores and biomechanical measurements.
Regression estimates a continuous numerical outcome, such as a post-treatment FMA-LE score.
Classification assigns a patient to a category, such as responder or non-responder, based on the
response definition used in the study. The selected paper evaluates univariate and multivariate linear
regression for continuous prediction, and Elastic Net and Artificial Neural Network models for
classification.\cite{bib1}
\section{Explainable Artificial Intelligence}
Explainable Artificial Intelligence (XAI) refers to methods that help people understand the reasoning
behind machine-learning predictions. In medical applications, clinicians need to understand which
features influence an estimated outcome.
The selected paper uses stepwise multivariate linear regression to select predictors associated with
post-treatment motor function and SHAP to estimate feature contributions to classification predictions.
These methods help connect model predictions with clinically meaningful characteristics, but they do
not prove causation.\cite{bib1}

\section{Important Clinical Assessment Measures}
Explainable Artificial Intelligence (XAI) refers to methods that help people understand the reasoning
behind machine-learning predictions. In medical applications, clinicians need to understand which
features influence an estimated outcome.
The selected paper uses stepwise multivariate linear regression to select predictors associated with
post-treatment motor function and SHAP to estimate feature contributions to classification predictions.
These methods help connect model predictions with clinically meaningful characteristics, but they do
not prove causation.\cite{bib1}
